\subsection{Global maximum on connected proper native schemes}

\begin{lemma}[Closed intermediates in finite-type curve chains]
\label{lem:fc-closed-plain-curve-chain}
\lean{EnvelopeProperty.FcSchemeCurveChains.closed_curveChain_of_connected_finiteType}
\leanok
A curve chain between two closed points of a finite-type scheme over an
algebraically closed field can be chosen with every intermediate point closed.
\end{lemma}
\begin{proof}
Induct on a given chain, proving the stronger assertion that its final point
may be replaced by any closed specialization. In the empty chain, closedness
of the initial point identifies the specialization with it. For the final
curve from the previous intermediate point to the endpoint, choose a closed
specialization of the previous intermediate point in its nonempty closed
singleton closure, using the Jacobson property. A closed curve containing a
point contains its singleton closure, so both new endpoints remain on the same
original integral curve. Apply induction to the preceding chain and append
this unchanged curve with the two new closed endpoints.
\end{proof}

\begin{lemma}[Maximum propagation at original closed supports]
\label{lem:fc-finite-type-closed-support-maximum}
\lean{EnvelopeProperty.Berkovich.RelativeSchemeValuedPoint.IsLocallyPSHOn.eq_of_isMaxOn_univ_of_closed_supports_finiteType}
\leanok
\uses{lem:fc-closed-plain-curve-chain,thm:fc-native-maximum-integral-closed-curve}
On a connected finite-type scheme over a complete algebraically closed
ultrametric normed field, an attained native locally psh global maximum at a
closed support has the same value at every native point with closed support.
\end{lemma}
\begin{proof}
Take the actual integral closed curve chain between the two original closed
supports, and choose closed intermediates by the preceding lemma. At each
intermediate construct a native point over the original field norm. Pullback
to each original integral curve propagates the attained maximum by the integral
scheme-curve maximum principle. The empty chain is handled by the integral
reduced singleton scheme. Induction through the finite chain gives equality.
\end{proof}

\begin{lemma}[Rational comparison on an original integral proper scheme]
\label{lem:fc-integral-proper-rational-comparison}
\lean{EnvelopeProperty.Berkovich.RelativeSchemeValuedPoint.IsLocallyPSHOn.eq_section_of_isMaxOn_univ_integral_proper,EnvelopeProperty.Berkovich.RelativeSchemeValuedPoint.IsLocallyPSHOn.eq_of_isMaxOn_section_integral_proper}
\leanok
\uses{lem:fc-finite-type-closed-support-maximum,lem:fc-rational-section-scalar-restriction}
Let an integral proper scheme over a complete algebraically closed ultrametric
normed field carry a native locally psh function. Any original native point
and the native point of any original rational section have the same value if
either point attains the original global maximum.
\end{lemma}
\begin{proof}
Use the completed algebraic closure of the arbitrary point's completed residue
field. The arbitrary point has an actual rational lift, and the original
rational section has its actual scalar-changed rational section. Their entire
native projections recover the original two points, and their lifted supports
are closed. The original integral proper scheme is geometrically connected:
its original global coefficient map is bijective by universal closedness and
algebraic closedness, and genuine flat field base change preserves that whole
field of global sections. Its actual scalar change is therefore connected.
The unchanged locally psh function pulls back to the scalar-changed scheme,
and the assumed original maximum pulls back to the appropriate lifted point.
Apply the preceding closed-support maximum statement, reversing its two
endpoints if needed, and use the exact native projection identities.
\end{proof}

\begin{theorem}[Connected proper native maximum principle]
\label{thm:fc-connected-proper-native-maximum}
\lean{EnvelopeProperty.Berkovich.RelativeSchemeValuedPoint.IsLocallyPSHOn.exists_closed_support_maximizer_of_isMaxOn_univ,EnvelopeProperty.Berkovich.RelativeSchemeValuedPoint.IsLocallyPSHOn.eq_of_isMaxOn_univ_connected_proper}
\leanok
\uses{lem:fc-integral-proper-rational-comparison,lem:fc-finite-type-closed-support-maximum,lem:fc-actual-rational-section}
On any connected proper scheme over a complete algebraically closed
ultrametric normed field, a native locally psh function attaining a global
maximum is constant.
\end{theorem}
\begin{proof}
First take the original support of a maximizing point. Its singleton closure,
with the vanishing-ideal reduced subscheme structure, is an actual integral
closed subscheme and remains proper over the original field. The whole native
point lifts to this closed subscheme. Choose an actual original rational
section there. The integral proper comparison shows that its image in the
original scheme attains the same maximum; this image has closed support.

For any other original native point repeat the singleton-closure construction
and choose an actual rational section of that integral proper subscheme. The
two section images have closed support in the original scheme, so connected
finite-type curve propagation gives them the same maximizing value. This
second section therefore maximizes the pulled-back function on its integral
proper subscheme. Apply integral proper comparison to the lifted arbitrary
point and project back. Its value equals the original attained maximum.
\end{proof}
